% !TEX root = ../general/main.tex
\section{Methodology}

\subsection{Causal framework}

We consider data from two sources: a randomised controlled trial (RCT) and a real-world data study (RWD).
The data-source indicator $S \in \{0, 1\}$ takes the value $S = 1$ for RCT observations and $S = 0$ for RWD observations.
The RCT contributes $n_1$ independent observations and the RWD contributes $n_0$ independent observations.
We observe a vector of pre-treatment covariates $X \in \mathbb{R}^p$ and a binary treatment $A \in \{0, 1\}$ for each individual.
The outcome of interest is the nonnegative survival time $T$, subject to censoring.
We allow two types of censoring: right censoring, and interval censoring.
For each subject we observe a pair $(L, R)$ with $0 \leq L \leq R \leq \infty$ satisfying $T \in [L, R]$, together with a censoring-type indicator $\delta \in \{0, 1, 2\}$.
We use the convention that $R = \infty$ stands for $T \in [L, \infty)$.
These correspond to exact observation ($\delta = 1$: $L = R = T$), right censoring ($\delta = 0$: $L = C$ and $R = \infty$, for some censoring time $C$), and interval censoring ($\delta = 2$: $0 \leq L < R < \infty$).
The standard right-censored setup is recovered when $\delta \in \{0, 1\}$ for every subject.
We write $\mathcal{C}$ for the inspection and censoring process underlying $(L, R, \delta)$, and refer to \citet{sun2006statistical} for background.
We adopt the potential-outcomes framework \citep{rubin1974}: $T(a)$, $\mathcal{C}(a)$, and the corresponding observations $(L(a), R(a), \delta(a))$ are defined analogously under treatment $A = a$.


We estimate the conditional average treatment effect (CATE):
\begin{equation}
\tau(x) = \E[\log T(1) - \log T(0) \mid X = x].
\label{eq:cate}
\end{equation}
The CATE captures the heterogeneity of the treatment effect across covariate values.
The corresponding estimand on the multiplicative time scale is the acceleration factor $\exp\{\tau(x)\}$.
The acceleration factor admits a direct interpretation on the survival-time scale: at every quantile of the survival distribution, treatment rescales the survival time by a constant multiplicative factor \citep{Brathovde2024}.
This time-scaling is more intuitive in a clinical setting than the relative change in event rate conveyed by a hazard ratio \citep{Swindell2009}.
The acceleration factor also enjoys causal properties that the hazard ratio lacks.
\citet{Brathovde2024} show that the observed acceleration factor identifies its causal counterpart under exchangeability and consistency.
This identification continues to hold under unmeasured frailty and treatment-effect heterogeneity.
The hazard ratio loses its causal interpretation in those same settings, because conditioning on survival induces a built-in selection bias \citep{hernan2010}.
The acceleration factor is collapsible: the marginal and conditional acceleration factors coincide when the treatment effect is homogeneous, regardless of the distribution of unmeasured baseline risk \citep{Aalen2015, Crowther2023}.

We impose conditions on the causal structure of the problem and on the censoring mechanism.
\vspace*{-\parskip}
\begin{adjustwidth}{1.5em}{1.5em}
\setlength{\parskip}{0pt}
\begin{assumption}[Consistency]\label{ass:sutva}
$T = A\,T(1) + (1 - A)\,T(0)$.
\end{assumption}

\begin{assumption}[Unconfoundedness of the RCT]\label{ass:rct-unconf}
$T(a) \,\ind\, A \mid X,\, S = 1$ for each $a\in\{0, 1\}$.
\end{assumption}

\begin{assumption}[Positivity]\label{ass:positivity-both}
$0 < \p(A = 1 \mid X = x,\, S = s) < 1$ for each $s \in \{0, 1\}$ and all $x$ in the support of $X \mid S = s$.
\end{assumption}

\begin{assumption}[Cross-source exchangeability]\label{ass:transport}
$T(a) \mid X,\, S = 1 \;\overset{d}{=}\; T(a) \mid X,\, S = 0$, for each $a \in \{0, 1\}$.
\end{assumption}

\begin{assumption}[Conditionally non-informative censoring]\label{ass:censoring}
$T(a) \,\ind\, \mathcal{C}(a) \mid X,\, A,\, S$ for each  $a\in\{0, 1\}$.
\end{assumption}

\end{adjustwidth}
\setlength{\parskip}{0.8em}
\vspace{0.8em}


\begin{remark}[Possible weakening to CATE-only transport]
Assumption~\ref{ass:transport} may be stronger than needed.
It could be replaced by the weaker condition that only the conditional log-treatment effect transports across sources:
\begin{equation}
\E[\log T(1) - \log T(0) \mid X,\, S = 1] \;=\; \E[\log T(1) - \log T(0) \mid X,\, S = 0].
\label{eq:cate-transport}
\end{equation}
This alleviation may suffice in our setting because the causal estimand is the CATE on the log-time scale.
\textit{To be verified.}
\end{remark}


Our identification strategy relaxes unconfoundedness in the RWD and retains transportability of the CATE across sources.
Assumption~\ref{ass:rct-unconf} holds by design in a randomised trial.
Treatment in the RWD may depend on confounders not contained in $X$.
The resulting bias is absorbed by the confounding function $c(x)$ introduced in the next section.
We assume the CATE to be time-invariant by not explicitly letting it depend on time.
Other approaches that combine evidence from multiple sources take the opposite route: they allow source-specific CATEs but require unconfoundedness within each source \citep{stuart2011propensity, dahabreh2019generalizing, shyr2025multistudy}.
Such heterogeneity in the CATE arises from varying distributions of unobserved effect modifiers \citep{degtiar2023review}, for example physical activity in breast-cancer treatment \citep{cannioto2021physical}.
The two relaxations cannot coexist in our setting.
Without unconfoundedness of the RWD, the CATE and $c(x)$ are not separately identifiable unless the CATE transports across sources.
We therefore retain Assumption~\ref{ass:transport} as the explicit cost of absorbing RWD confounding into $c(x)$, and acknowledge that CATE heterogeneity across sources falls outside the scope of the current framework.

\subsection{Accelerated failure time decomposition}

The confounding function captures the bias in the RWD that the causal effect $\tau$ alone cannot explain.
Unconfoundedness fails in the RWD by design.
The treated-versus-control contrast observed in the RWD therefore differs from the CATE in general.
We define the \emph{confounding function} $c(x)$ as:
\begin{equation}
c(x) \;=\; \E[\log T \mid X = x,\, A = 1,\, S = 0] \,-\, \E[\log T \mid X = x,\, A = 0,\, S = 0] \,-\, \tau(x).
\label{eq:cf}
\end{equation}
The confounding function represents the difference between the RWD treated-versus-control contrast and the true causal effect.
Under Assumptions~\ref{ass:sutva}--\ref{ass:transport}, the conditional log survival time admits an additive decomposition into a baseline, the causal effect, and the confounding bias.
\begin{proposition}[Accelerated failure time decomposition]\label{prop:decomp}
Under Assumptions~\ref{ass:sutva}--\ref{ass:transport}, the conditional log survival time satisfies:
\begin{equation}
\E[\log T \mid A,\, X,\, S] \;=\; m_0(X, S) \,+\, \tau(X)\,A \,+\, (1 - S)\,A\,c(X),
\label{eq:decomp}
\end{equation}
where $m_0(X, S) := \E[\log T \mid A = 0,\, X,\, S]$ is the baseline log survival time.
\end{proposition}

The decomposition separates three pieces.
The baseline $m_0(X, S)$ is the control-arm log survival time in each source.
The CATE $\tau(X)$ enters additively whenever an observation is treated, regardless of source.
The confounding function $c(X)$ contributes only in the treated arm of the RWD, reflecting that the RWD treated-versus-control contrast carries bias on top of the causal effect.

\begin{proposition}[Identification]\label{prop:ident}
Under Assumptions~\ref{ass:sutva}--\ref{ass:transport}:
\begin{enumerate}[label=\textup{(\roman*)}]
\item The CATE is identified from the RCT alone:
\begin{equation}
\tau(x) \;=\; \E[\log T \mid X = x,\, A = 1,\, S = 1] \,-\, \E[\log T \mid X = x,\, A = 0,\, S = 1].
\label{eq:ident-tau}
\end{equation}
\item Given $\tau$, the confounding function is identified from the RWD.
\item Neither $\tau$ nor $c$ is identified from the RWD alone; the RWD identifies only the composite $\tau(x) + c(x)$.
\end{enumerate}
\end{proposition}

Proofs of Propositions~\ref{prop:decomp} and~\ref{prop:ident} are in Supplementary Materials~\ref{sm:proofs}.
The same identification result holds for the acceleration factor $\exp\{\tau(x)\}$ by continuity of the exponential.

We illustrate the confounding function through a structural causal model with an unmeasured confounder $U$ as a special case in Supplementary Materials~\ref{sm:scm}.
This reading gives $c(x)$ a mechanistic interpretation.
It decomposes $c(x)$ into a prognostic imbalance between the RWD treatment arms and effect modification by $U$.
The decomposition motivates the prior placed on $c$.

% Joint estimation across both sources improves the efficiency of estimating $\tau$, even though the CATE is identified from the RCT alone.
% The decomposition lets the RWD contribute through the shared baseline $m_0(X, S)$ and through the additional treated observations, while the confounding function $c$ absorbs the bias.
% The efficiency gain is largest when $c$ is small or smoothly varying, so that few degrees of freedom are spent on the bias correction.

We model the log survival time with the accelerated failure time (AFT) specification implied by the decomposition:
\begin{equation}
\log T_i \;=\; m_0(X_i, S_i) \,+\, A_i\,\tau(X_i) \,+\, (1 - S_i)\,A_i\,c(X_i) \,+\, \varepsilon_i.
\label{eq:aft-model}
\end{equation}
We model the errors $\varepsilon_i$ as independent with mean zero and a nonparametric distribution.
The regression functions $m_0$, $\tau$ and $c$ retain the causal meanings established above.
We assign each a separate Bayesian tree ensemble prior.


\subsection{Nonparametric error distribution}

We model the errors $\varepsilon_i$ in \eqref{eq:aft-model} with a hierarchical Dirichlet process mixture of normals (HDPM).
The prior is source-specific: $\varepsilon_i \mid S_i = s \sim F_s$.
The model lets $F_0$ and $F_1$ differ flexibly while sharing information about which residual types are present across sources.
Existing data-fusion methods typically assume $F_0 = F_1$, for example via a single Gaussian likelihood.
This assumption is unsatisfactory in our setting.
Randomised trials and real-world data sources typically differ in measurement protocol, follow-up quality, and the prevalence of extreme outcomes.
These differences plausibly enter the residual distribution rather than the structural mean.
The two error distributions are unlikely to be entirely unrelated.
A model that treats them as fully independent forfeits useful pooling, particularly when one source is small.

Conditional on the source indicator $S_i = s$, we model the residual as a location mixture of Gaussians with source-specific scale $\sigma_s > 0$ and mixing measure $G_s$:
\begin{equation}
\varepsilon_i \mid S_i = s,\, G_s,\, \sigma_s \;\sim\; \int \frac{1}{\sigma_s}\, \phi\!\left(\frac{w - \theta}{\sigma_s}\right) dG_s(\theta).
\label{eq:hdp-fs}
\end{equation}
Gaussian-mixture priors of this form have been used for AFT residuals in single-source BART models \citep{henderson2020individualized}.
The mixing measures share latent structure through the hierarchical Dirichlet process of \citet{teh2006hierarchical}, with a top-level random measure $G^*$:
\begin{equation}
G^* \;\sim\; \mathrm{DP}(\gamma, H), \qquad G_s^* \mid G^*,\, M_s \;\sim\; \mathrm{DP}(M_s, G^*), \qquad s \in \{0, 1\},
\label{eq:hdp-prior}
\end{equation}
with base distribution $H = \mathcal{N}(0, \sigma_\theta^2)$.
By Sethuraman's stick-breaking representation \citep{sethuraman1994constructive}, the top-level measure is discrete:
\begin{equation}
G^* \;=\; \sum_{k=1}^\infty \omega_k\, \delta_{\theta_k^*}, \qquad \theta_k^* \stackrel{\mathrm{iid}}{\sim} H, \qquad \omega_k = u_k \prod_{l < k}(1 - u_l), \qquad u_k \stackrel{\mathrm{iid}}{\sim} \mathrm{Beta}(1, \gamma).
\label{eq:hdp-stick}
\end{equation}
Each source-specific measure $G_s^*$ is supported on the same atoms $\{\theta_k^*\}$ as $G^*$, with its own weight vector $\bm{\pi}_s = (\pi_{sk})_{k \ge 1}$.
The shape of each $F_s$ is therefore parameterised by a shared set of residual types and a source-specific weighting over them.

For $m_0$, $\tau$, and $c$ to retain their interpretation as the conditional-mean components defined in \eqref{eq:decomp}, each $F_s$ must have mean zero.
We follow \citet{yang2010semiparametric} and obtain centred mixing measures by subtracting the source-specific mean of the unconstrained measure:
\begin{equation}
\mu_s \;=\; \int \theta\, dG_s^*(\theta) \;=\; \sum_{k=1}^\infty \pi_{sk}\, \theta_k^*, \qquad G_s \;=\; \sum_{k=1}^\infty \pi_{sk}\, \delta_{\theta_k^* - \mu_s}.
\label{eq:hdp-centring}
\end{equation}
By construction $\E[\varepsilon_i \mid S_i = s,\, G_s,\, \sigma_s] = 0$.
The unconstrained atoms $\theta_k^*$ remain shared across sources, while the centred atoms $\theta_k^* - \mu_s$ are source-specific through the shifts $\mu_0$ and $\mu_1$.
We place a shared scaled-inverse-$\chi^2$ prior $\sigma_s^2 \sim \nu\lambda / \chi^2_\nu$ on the per-source residual scales, with degrees of freedom $\nu$ and scale $\lambda$.
We set $\nu = 3$ following the BART default and calibrate $\lambda$ to the empirical residual variance following \citet{chipman2010bart}.

\subsection{Overview of BART}

Bayesian additive regression trees \citep{chipman2010bart} provide a flexible nonparametric prior over an unknown function $f : \mathbb{R}^p \to \mathbb{R}$.
The prior represents $f$ as a sum of regression trees:
\begin{equation}
f(X) \;=\;
  \sum_{j=1}^{J} g(X;\, \mathcal{T}_j, \mathcal{H}_j),
\label{eq:bart-revised}
\end{equation}
where each $\mathcal{T}_j$ is a binary tree with $b_j$ terminal nodes.
Each interior splitting rule has the form $\{x_\ell \leq x^*\}$ for some covariate index $\ell$ and splitting value $x^*$.
We write $\mathcal{H}_j = \{h_{j,1}, \ldots, h_{j,b_j}\}$ for the step heights at those terminal nodes.
The function $g(X;\, \mathcal{T}_j, \mathcal{H}_j)$ routes $X$ through the splits of $\mathcal{T}_j$ to a terminal node $r$ and returns the step height $h_{j,r}$ at that node.

The prior on $(\mathcal{T}_1, \mathcal{H}_1), \ldots, (\mathcal{T}_J, \mathcal{H}_J)$ is independent across trees.
It decomposes into a tree-structure prior and a step-height prior given the structure.
The structure of each $\mathcal{T}_j$ follows the recursive splitting process of \citet{chipman1998bayesian}.
A node at depth $d$ is non-terminal with probability $\alpha(1+d)^{-\beta}$.
This generative process is a heterogeneous Galton--Watson branching process, in which the per-generation offspring distribution depends on the depth $d$ \citep{rockova2019theory}.
The hyperparameters $\alpha$ and $\beta$ regularise tree depth.
Conditional on the tree, the splitting variable $x_\ell$ at each interior node is drawn uniformly from the available covariates.
The splitting value $x^*$ is drawn uniformly from its observed values.
The step heights receive independent conjugate normal priors $h_{j,r} \mid \mathcal{T}_j \sim \mathcal{N}(0,\, \sigma_h^2)$, parameterised as $\sigma_h = k/(2\sqrt{J})$ with $k > 0$ the step-height scale.
\citet{chipman2010bart} centre and rescale the response, then choose $k$ so that the implied prior on $f$ assigns circa $95\%$ of its mass to the observed response range.
We denote the BART prior with these hyperparameters by $\mathrm{BART}(J, k, \alpha, \beta)$.
The four arguments are the number of trees $J$, the step-height scale $k$, and the tree-structure hyperparameters $\alpha$ and $\beta$.


\subsection{Prior specification}

We place independent Bayesian tree priors on the components of the decomposition~\eqref{eq:decomp}.
We tune the regularisation of each component through the tree-structure prior and the ensemble size to match the substantive role of each component.

We allow for heterogeneity in the baseline effect between sources.
The prognostic function $m_0$ is split into a shared component and a source-specific deviation so that the prior centres on between-source borrowing.
We write:
\begin{equation}
m_0(X, S) \;=\; m_0^{\mathrm{sh}}(X) + (1 - S)\, d(X),
\label{eq:m0-decomp-prop}
\end{equation}
where $m_0^{\mathrm{sh}}$ is shared across sources and $d$ is active only in the RWD.
The decomposition writes the difference between the two source-specific prognostic functions as $d(X)$.
The prior on $d$ has mean zero, so the prior centres on equality of the two prognostic functions.
An alternative would place a single BART on $m_0(X, S)$ with $S$ as an ordinary covariate.
Such a model does not encode this prior centring.
The prior on $d$ then provides a direct dial for the strength of borrowing.
We use the tree-structure hyperparameters $(\alpha_{\mathrm{sh}}, \beta_{\mathrm{sh}}) = (0.95, 2)$ and $J_{\mathrm{sh}} = 200$ trees for $m_0^{\mathrm{sh}}$.
The large ensemble lets $m_0^{\mathrm{sh}}$ accommodate complex disease-intrinsic structure.
We use the same hyperparameters $(\alpha_d, \beta_d) = (0.95, 2)$ for $d$ but only $J_d = 50$ trees.
The smaller ensemble suffices because $d$ only needs to absorb residual differences in protocol and patient population.
This construction coincides with a heterogeneous meta-analytic-predictive prior on the prognostic function (Supplementary Materials~\ref{sm:map-m0}).

We place a depth-penalised BART prior on the treatment-effect function $\tau$ to discourage higher-order interactions in $X$.
We set the depth parameter $\beta_\tau = 3$ and $\alpha_\tau = 0.95$.
The prior mass concentrates on trees of depth one or two.
The number of trees is reduced relative to the baseline prognosis to $J_\tau = 100$ trees.
This reflects the expectation that the treatment-effect surface is simpler than the baseline prognosis.

We place a more strongly regularised BART prior on the confounding function $c$.
We set the base splitting probability $\alpha_c = 0.25$ and the depth parameter $\beta_c = 3$.
We set the number of trees to $J_c = 50$.
This discourages splits at the root and deep interactions when splits do occur.


\begin{remark}[Multiplicative shrinkage on $c$]
An alternative representation is $c(x) = \kappa\, g_c(x)$, with $\kappa$ a scalar shrinkage parameter and $g_c \sim \mathrm{BART}$.
A heavy-tailed prior such as $\kappa \sim \mathcal{C}^+(0, 1)$ concentrates the posterior at small magnitudes while permitting occasional large values.
A normal prior $\kappa \sim \mathcal{N}(0, \sigma_\kappa^2)$ is a softer alternative.
Either choice imposes stronger regularisation on $c(x)$ than the sole BART prior, and the posterior of $\kappa$ admits a reading as a level of RWD confounding.
\end{remark}

We place a shared scaled-inverse-$\chi^2$ prior $\sigma_s^2 \sim \nu\lambda / \chi^2_\nu$ on the per-source residual scales with $\nu$ degrees of freedom and scale $\lambda$.
We set $\nu = 3$ following the BART default and calibrate $\lambda$ to the empirical residual variance following \citet{chipman2010bart}.
We standardise the log survival time before fitting by the centring and scaling constants of a preliminary log-normal AFT fit.
This rescales the response to a unit-variance scale.
Each forest carries its own prior step-height scale $k_f$ with $f \in \{m_0^{\mathrm{sh}}, d, \tau, c\}$.
The step-height scales calibrate the prior magnitude of each ensemble to match the scale of the standardised response.
We recommend $k_f = 1$ for all four forests as a default and allow further tuning via cross-validation.

We summarise the full hierarchical model:
\begin{align*}
\log T_i &\;=\; m_0^{\mathrm{sh}}(X_i) + (1 - S_i)\, d(X_i) + A_i\, \tau(X_i) + (1 - S_i)\, A_i\, c(X_i) + \varepsilon_i, \\
m_0^{\mathrm{sh}} &\;\sim\; \mathrm{BART}(200,\, k_{\mathrm{sh}},\, 0.95,\, 2), \\
d &\;\sim\; \mathrm{BART}(50,\, k_d,\, 0.95,\, 2), \\
\tau &\;\sim\; \mathrm{BART}(100,\, k_\tau,\, 0.95,\, 3), \\
c &\;\sim\; \mathrm{BART}(50,\, k_c,\, 0.25,\, 3), \\
\varepsilon_i \mid S_i = s,\, G_s,\, \sigma_s &\;\sim\; \int \sigma_s^{-1}\, \phi\!\left(\tfrac{w - \theta}{\sigma_s}\right) dG_s(\theta), \\
G_s &\;=\; \textstyle\sum_k \pi_{sk}\, \delta_{\theta_k^* - \mu_s}, \quad s \in \{0, 1\}, \\
G_s^* \mid G^*,\, M_s &\;\sim\; \mathrm{DP}(M_s,\, G^*), \quad s \in \{0, 1\}, \\
G^* &\;\sim\; \mathrm{DP}(\gamma,\, H), \qquad H = \mathcal{N}(0,\, \sigma_\theta^2), \\
\sigma_s^2 &\;\sim\; \nu\lambda / \chi^2_\nu, \quad s \in \{0, 1\}.
\end{align*}

We refer to this combined model as the \emph{Bayesian fusion forest}.






\subsection{Posterior inference*}

\textit{This section is largely unfinished.
The paragraphs on average treatment effects and posterior computation are ready.
The following items are ideas for material that could go in this section.}

\begin{itemize}
  \item \textbf{Treatment-effect estimands.}
    Posterior summaries of the CATE $\tau(x)$ and the acceleration factor $\exp\{\tau(x)\}$ at fixed $x$: posterior mean and credible bands.
    Investigate the impact of covariates on the outcome, e.g. via linear projections or TE-VIMs.
  \item \textbf{Heterogeneity.}
    Posterior measures of heterogeneity of the treatment effect.
  \item \textbf{Borrowing.}
    Posterior measures of borrowing: behaviour of the deviation forest $d$ relative to the shared baseline $m_0^{\mathrm{sh}}$, magnitude of the borrowing scale $\sigma_d$, posterior mass on $c \equiv 0$ et cetera.
  \item \textbf{Residual fit.}
    Was the heavy nonparametric model for the error distribution justified, or would a simpler model suffice?
    Posterior number of active mixture clusters and between-source comparison of $F_0$ and $F_1$ for example.
\end{itemize}


We compute average treatment effects by marginalising the conditional effect $\tau(\cdot)$ against a target covariate distribution.
We take the target to be the union of the two source studies, with covariate distribution $F_X = \pi_0 F_X^0 + \pi_1 F_X^1$.
Here $F_X^s$ is the covariate distribution within source $s \in \{0, 1\}$ and $\pi = (\pi_0, \pi_1)$ lies on the unit simplex.
The posterior of the marginal effect carries uncertainty from three sources: the posterior over $\tau$, the within-source covariate distributions $F_X^s$, and the mixing fractions $\pi$.
We propagate all three jointly by a hierarchical Bayesian bootstrap.
The choice of $\pi$ encodes the analyst's target population.
Natural deterministic choices include the empirical fractions $\pi_s = n_s / n$ and equal weighting $\pi = (1/2, 1/2)$.
The source-specific limits $\pi = (1, 0)$ and $\pi = (0, 1)$ target the population of the RWD and RCT respectively.
We may instead place a Dirichlet prior $\pi \sim \mathrm{Dirichlet}(\alpha)$ on the mixing fractions, with concentration parameter $\alpha = (\alpha_0, \alpha_1)$.
This propagates uncertainty about the target composition.
The choice $\alpha = (n_0, n_1)$ centres the prior on the empirical fractions and approximately recovers a pooled Bayesian bootstrap.
We recommend this as the default.
The choice $\alpha = (1, 1)$ gives a flat prior on the simplex.
We draw mixing fractions $\pi^{(b)} \sim \mathrm{Dirichlet}(\alpha)$ for each posterior draw $\tau^{(b)}(\cdot)$ of the conditional effect.
We draw within-source weights $w^{(s, b)} \sim \mathrm{Dirichlet}(\mathbf{1}_{n_s})$ over the observations of source $s$ independently for $s \in \{0, 1\}$.
The corresponding draw of the average treatment effect is:
\begin{equation}
\bar{\tau}^{(b)} \;=\; \sum_{s \in \{0, 1\}} \pi_s^{(b)} \sum_{i: S_i = s} w_i^{(s, b)}\, \tau^{(b)}(X_i).
\label{eq:ate-bb}
\end{equation}
The within-source weights are the Bayesian bootstrap of \citet{rubin1981} applied to each $F_s$.
 

We sample from the posterior via a blocked Gibbs sampler built on an extended Bayesian backfitting approach \citep{hastie2000bayesian}.
We update the BART forests $m_0^{\mathrm{sh}}$, $d$, $\tau$, and $c$ in turn against the partial residual.
The partial residual subtracts the other forests and the current draw of the residual mixture from the augmented outcome.
Each update reduces to a standard BART update on a Gaussian working response \citep{chipman2010bart}.
We augment the censored observations \citep{tanner1987}.
The HDPM error block is updated by a separate Gibbs step over its cluster assignments, atoms, mixture weights, and concentration parameters.
We give the full sampler in the Supplementary Materials.




